\subsection{Causal sets}\label{sec:causets}

Let $\mathsf{CS}$ have locally finite posets as objects and order-embeddings
with order-convex image as morphisms. A becoming theory on $\mathsf{CS}$ is
defined as before, with (CB) read as $x\le b\Rightarrow x\,B\,b$.

\begin{proposition}\label{prop:causets}
Every becoming theory on $\mathsf{CS}$ is either the causal order in every
causal set or the universal relation in every causal set.
\end{proposition}

\begin{proof}
If $x,b$ are incomparable, $\{x,b\}$ is convex, so both embeddings of the
two-element antichain are morphisms: the verdict on incomparable pairs is a
single symmetric bit $\beta$, the same in every causal set. If $x>b$, the
interval $I=[b,x]$ is finite and convex, so $x\,B\,b$ depends only on $I$. If
$\beta=1$, in $I\sqcup\{z\}$ with $z$ incomparable to all, $x\,B\,z\,B\,b$; so
every such verdict is positive and $B$ is universal. If $\beta=0$ and
$x\,B\,b$ for some $x>b$, adjoin $e$ with $e<x$ only: $I$ remains convex, so
$x\,B\,b$ there, $e\,B\,x$ by (CB), and $e\,B\,b$ with $e$ incomparable to $b$,
contradicting $\beta=0$.
\end{proof}

Wüthrich and Callender state that ``The analogue of this theorem, however, is
straightforwardly false in CST'' \citep[\S3]{WC2017}. Their counterexample, on
the diamond causal set $d<b,c<a$, adds the incomparable pair $\{b,c\}$ to the
causal order and is invariant under that causal set's automorphisms. It is not
the restriction of any becoming theory in the present sense: relating $b$ to
$c$ fixes $\beta=1$, and Proposition~\ref{prop:causets} then relates $b$ to $d$,
which their relation excludes. Their own footnote~5 observes that the
counterexample ``could be a small proper part of a much larger causet''; (LC) is
exactly the requirement that the relation on such a part agree with the
relation in the larger causal set. The diagnosis matches the continuum case:
invariance within one structure is too weak, and locality across structures
restores Stein's theorem, here with no residual freedom because a causal set
has no light-cone boundary. Their claim is true on its own terms; the present
result uses a stronger, cross-structure hypothesis. Note also that
$\mathsf{CS}$ includes posets that are not manifoldlike (antichains, disjoint
unions), and the proof uses them: this is the causal-set counterpart of the
off-shell premise of Remark~\ref{rem:offshell}.

\subsection{Gödel universes are isolated}\label{sec:godel}

\begin{proposition}\label{prop:godel}
Extend $\GH$ to all oriented, time-oriented spacetimes with the same
morphisms. If $M$ is totally vicious ($p\in\Ichr_{\mathrm s}(p)$ for all $p$),
then every morphism into $M$ is onto, every theory satisfying (CB) has
$B_M=\U_M$, and there is no morphism between $M$ and any causal spacetime, in
either direction.
\end{proposition}

\begin{proof}
In a totally vicious spacetime $\Ichr(p)=\Ifut(p)=M$ for every $p$ (MS \S3.1).
Let $V$ be nonempty, open and causally convex, $p\in V$ and $q\in M$. There are
timelike curves $p\to q$ and $q\to p$; their concatenation is a causal curve
with both endpoints in $V$, so it lies in $V$, and $q\in V$. Thus $V=M$
(compare MS Remark~2.12). (CB) gives $B_M\supseteq\Ichr_M=M\times M$. A morphism
$N\to M$ from a causal $N$ would be onto, hence an isometry, which is
impossible since $M$ has closed timelike curves; a morphism $M\to N$ would map
a closed timelike curve of $M$ to one of $N$.
\end{proof}

Gödel's universe is totally vicious, so universal becoming there is forced;
this is Gödel's first step, in Stein's terms. But local covariance transmits
constraints only along morphisms, and by Proposition~\ref{prop:godel} no
morphism connects a totally vicious spacetime with a globally hyperbolic one. Gödel's modal step, from what holds
in such a universe to what holds in ours \citep[pp.~561--562]{Godel1949}, as
quoted in \citep[\S6]{Monton2006}, requires that becoming be the same
feature across nomologically possible universes. Local covariance is one
precise way of requiring a rule to be the same across universes, and it
transmits nothing out of totally vicious ones. A version of Gödel's step that does go through
would need a different cross-universe link. Proposition~\ref{prop:nonlocal} shows what such links
look like: natural but non-local. Compare the conceptual responses of
\citet{Savitt2005}, \citet{Dieks2006} and \citet{Aames2022}, which locate
becoming locally along worldlines.
